% Integration candidate; not included in main.tex until proof/evidence review.
\subsection{Observable calibration and unknown initialization}
A covariance probe must be computable from the learner's observations.
If its raw output is $Z=\boldsymbol1\vartheta+W_j$, subtracting the unknown target $\vartheta$ would reveal unavailable information.
Instead choose the first-difference matrix $D$ and use the observable contrast $Y=DZ=DW_j$.
This classical error-contrast construction~\cite{harville1974contrasts} gives covariance $K_j=D\Sigma_jD^\top$ and exact Gaussian affinity
\[
\rho_D=\det(K_0)^{1/4}\det(K_1)^{1/4}/\det[(K_0+K_1)/2]^{1/2}.
\]
Each directional error of its zero-threshold likelihood rule is at most $\rho_D$.
The common state at probe exit has conditional mean $m_j(Y)$ and explained variance $s_j^2=\operatorname{Var}_j(m_j(Y))$, computed from this projected observation law.
The supplement derives a linear-time two-filter implementation and the explicit finite-dimensional information cost of removing an unknown target.
Every discarded probe observation still consumes both budgets.

Let the actual initial error be an unknown fixed $e\in[-r,r]$, with public $r>0$.
Fix a regime prior $\pi_j>0$, mixture law $g$, and model posterior $p_j(y)$.
For each frozen post-probe program let $J_{j,a}^{r}=d_a^2r^2+\operatorname{Var}_j(\sum_t\omega_{a,t}\bar X_t)$, including its exact remaining horizon.
Its even conditional score is
\begin{equation}
L_a(y)=\sum_jp_j(y)\{J_{j,a}^{r}+k_a^2[m_j(y)^2-s_j^2]\}.
\label{eq:conditional-even-score}
\end{equation}
This score uses neither the target nor the sign of the initial error.
Write $F_a(y)=\sum_jp_j(y)J_{j,a}^{r}$, let $u$ minimize $L_a$ and $c$ minimize $F_a$, and set $S(y)=\sum_jp_j(y)[m_j(y)^2-s_j^2]$.
Both selectors use a fixed even tie rule.
For an even selector write $R_\pi(v)=\sup_{|e|\le r}R_{\pi,e}(v)$, and let $\varphi$ denote the standard normal density.

\begin{theorem}[Conditional value and its state-information ceiling]
For a fixed paid probe and finite scalar continuing-program catalogue with common persistence and the same deterministic learner weights under each regime,
\begin{equation}
\inf_v\sup_{|e|\le r}R_{\pi,e}(v)=\E_g\min_aL_a(Y),
\label{eq:conditional-minimax-value}
\end{equation}
where $v$ ranges over all randomized $Y$-measurable selectors in the invariant-statistic information class.
The selector $u$ attains the value, while its improvement over the posterior-marginal rule satisfies
\begin{align}
0\le R_\pi(c)-R_\pi(u)&\le K,\nonumber\\
K&=4\varphi(1)\bigl(\max_ak_a^2-\min_ak_a^2\bigr)\sum_j\pi_js_j^2.
\label{eq:state-selection-ceiling-main}
\end{align}
For the running mean and a stable averaged affine-delay learner with affordable horizons proportional to budget scale $B$, $K=O(B^{-2})$.
\label{thm:conditional-value}
\end{theorem}

Conditioning and pointwise loss minimization are classical decision tools~\cite{moustakides2012joint,vasconcelos2020observation,reinhard2026minimax}.
Here symmetrizing a selector cancels its unknown-initialization cross term, and averaging its risks at $e=\pm r$ proves the attainable bound in Eq.~\eqref{eq:conditional-minimax-value}.
The ceiling follows from $L_c-L_u\le(\max k_a^2-\min k_a^2)|S|$ and $\E_g|S|\le4\varphi(1)\sum_j\pi_js_j^2$.
The supplementary proofs give the stable-learner coefficient bounds and the vector extension, whose covariance dual uses classical minimax interchange~\cite{sion1958minimax} and may require tie randomization.

\subsection{When identification can outperform all fixed prefixes}
Theorem~\ref{thm:conditional-value} separates knowledge of the covariance regime from knowledge of the current noise state.
The latter has a second-order value for averaged learners; the former can change the leading risk coefficient.
For a running mean with budgets $(\beta_mB,\beta_eB)$, define
\[
a_a=\min\{\beta_m/(h+q_a),\beta_e/b_a\},\quad
K_j(a)=\frac{q_a^{-1}+\theta_j(1+\lambda^{b_a})/(1-\lambda^{b_a})}{a_a}.
\]
Let $a_j^*$ minimize $K_j$, $K_\pi=\min_a\sum_j\pi_jK_j(a)$,
$\Delta=K_\pi-\sum_j\pi_jK_j(a_j^*)$, and
$A=\sum_j\pi_j[K_j(a_{1-j}^*)-K_j(a_j^*)]$.
The full-budget fixed comparison permits every affordable training prefix and stop, at the same learner and step.

\begin{corollary}[First-order separation and a no-gain boundary]
For a finite catalogue, persistence bounded below one and bounded delay, a paid target-invariant probe with fixed directional error bound $\delta$ certifies superiority to every fixed prefix and fixed random mixture for sufficiently large $B$ if $\Delta>\delta A$.
The certified gain is at least $(\Delta-\delta A)/B-C(n+1)/B^2$, where $n$ is probe length and $C$ is finite.
If the regimes share a leading minimizing action, no selector from the maximal-horizon post-probe catalogue plus stop following a length-$O(\log B)$ probe gains more than $O(\log B/B^2)$ over the best fixed prefix.
\label{cor:leading-prefix-separation}
\end{corollary}

The supplement proves the exact-prefix expansion and gives a rational witness on the public grid with a positive leading margin of approximately $0.07765$ at directional bound $0.025$.
Finite-budget admission still uses the actual projected probe, exact horizons and signed risk bounds, rather than this asymptotic margin.
In the Gaussian implementation, probing is admitted only when its prior-weighted upper risk is below the best full-budget fixed-prefix risk; otherwise the controller executes that fixed prefix without probing.
This is a resource-program comparison at a declared learner and model, rather than dominance over unrestricted estimators.
